\chapter{DESENVOLVIMENTO DO MODELO DE NEGÓCIO}
\label{cap:elementos}

Este capítulo consolida o framework adotado pelo Grupo 3 para estruturar \textit{startups} de base tecnológica. A proposta integra cinco dimensões que, em conjunto, aumentam a probabilidade de construir um modelo de negócio robusto, escalável e financeiramente sustentável.

Em perspectiva executiva, o capítulo demonstra como decisões estratégicas, desenho organizacional e disciplina de execução convergem para transformar intenção empreendedora em resultado mensurável.

\section{Framework Integrado de Cinco Dimensões}

\subsection{Dimensão estratégica: PMBOK 7 e tailoring dinâmico}

Na dimensão estratégica, o foco está na geração de valor contínuo e no aprendizado validado. Em vez de processos rígidos, adota-se uma lógica orientada a princípios, com três eixos principais:

\begin{itemize}
    \item \textbf{Stewardship e governança participativa}: decisões orientadas ao valor entregue ao cliente e ao negócio;
    \item \textbf{Tailoring por maturidade}: adaptação dos ritos de gestão ao estágio da \textit{startup};
    \item \textbf{ROI mensurável}: cada ciclo de execução precisa produzir evidência de valor ou hipótese validada.
\end{itemize}

Como desdobramento operacional, essa dimensão se materializa em três frentes. Primeiro, define-se um ciclo curto de planejamento e revisão (quinzenal ou mensal), com hipóteses explícitas de mercado, produto e monetização. Em seguida, estabelece-se um quadro de priorização que classifica iniciativas por impacto esperado, esforço e risco de execução. Por fim, consolida-se uma rotina de \textit{review} estratégica, na qual os resultados do período são comparados com os objetivos definidos.

Para manter coerência entre estratégia e execução, recomenda-se monitorar: taxa de hipóteses validadas por ciclo, tempo médio entre decisão e experimento, e retorno incremental por iniciativa priorizada. Esses indicadores permitem ajustar rapidamente o direcionamento da \textit{startup} sem perder foco no valor entregue.

\subsection{Dimensão estrutural: platform design}

O modelo de negócio é concebido como um ecossistema multilateral, e não como cadeia linear de produção. Essa escolha amplia efeitos de rede e potencial de escala por meio de três motores:

\begin{itemize}
    \item \textbf{Motor de transações}: redução de fricção entre os participantes da plataforma;
    \item \textbf{Motor de aprendizado}: acumulação de dados e aumento de defensibilidade competitiva;
    \item \textbf{Feedback loops}: ciclos rápidos de retroalimentação entre produto, usuários e mercado.
\end{itemize}

Na prática, os desdobramentos estruturais envolvem mapear com precisão os atores do ecossistema (usuários finais, parceiros, fornecedores de tecnologia e canais de distribuição), definir a proposta de valor específica para cada lado da plataforma e estabelecer regras de interação que reduzam atrito de entrada. Esse desenho deve considerar mecanismos de atração inicial (\textit{seeding}), incentivos de permanência e barreiras de saída para aumentar retenção.

Também é necessário definir a lógica de monetização da plataforma desde cedo, mesmo que em versão preliminar. Entre os formatos mais usuais estão comissão por transação, assinaturas, \textit{tiers} de funcionalidades e serviços complementares. A escolha deve ser testada com indicadores como crescimento de base ativa, taxa de conversão e receita média por usuário.

\subsection{Dimensão de engenharia: Flow Framework}

A operação técnica é conectada diretamente ao desempenho de negócio. Para isso, o fluxo de trabalho é monitorado em quatro classes:

\begin{itemize}
    \item \textbf{Features}: entregas que geram valor percebido e potencial de receita;
    \item \textbf{Defects}: falhas que reduzem qualidade e experiência do usuário;
    \item \textbf{Riscos}: exposição a vulnerabilidades e não conformidades;
    \item \textbf{Dívida técnica}: parcela de capacidade reservada à sustentabilidade do código.
\end{itemize}

Recomenda-se reservar aproximadamente 20\% da capacidade da equipe para manutenção evolutiva e refatoração, mitigando crescimento desordenado da dívida técnica.

Como desdobramento de gestão técnica, cada classe de trabalho deve possuir política de entrada e critério de qualidade de saída. Em \textit{Features}, prioriza-se entrega incremental com validação de uso real; em \textit{Defects}, resposta proporcional ao impacto no cliente; em \textit{Riscos}, mitigação preventiva com monitoramento contínuo; e em dívida técnica, execução planejada para preservar velocidade futura de entrega.

Para conectar engenharia e negócio, recomenda-se acompanhar \textit{lead time}, \textit{throughput}, taxa de retrabalho e disponibilidade do serviço. A interpretação conjunta dessas métricas reduz decisões baseadas apenas em percepção e fortalece a previsibilidade da operação.

\subsection{Dimensão fiscal: Lei do Bem como alavanca de receita indireta}

A dimensão fiscal transforma atividades de pesquisa, desenvolvimento e inovação (PD\&I) em vantagem econômica por meio de conformidade estruturada. Com rastreabilidade adequada de iniciativas técnicas, torna-se possível recuperar parcela relevante dos custos de engenharia, com potencial de até 34\% conforme o enquadramento legal.

Os principais desdobramentos dessa dimensão são organizacionais e documentais. Do ponto de vista organizacional, é necessário delimitar quais frentes de trabalho se enquadram como PD\&I, com objetivos técnicos, incertezas e resultados esperados claramente registrados. Do ponto de vista documental, exige-se rastrear horas, artefatos de desenvolvimento, evidências de experimentação e resultados obtidos em cada ciclo.

Quando a governança fiscal é incorporada à rotina do time, o benefício não se limita ao incentivo tributário: há melhoria na disciplina de execução, maior clareza sobre custos de inovação e aumento da capacidade de justificar investimentos técnicos perante sócios e potenciais investidores.

\subsection{Governança de valor: VMO e OKRs}

A governança proposta substitui uma lógica centrada em controle burocrático por gestão orientada a valor. Assim, o \textit{Value Management Office} (VMO) organiza a priorização com uso de OKRs e OrKRs (\textit{Objectives, Risks and Key Results}), articulando estratégia, risco e mensuração contínua.

Como desdobramento prático, o VMO atua como camada de integração entre liderança, produto e engenharia. Sua função é tornar explícita a cadeia de valor de cada iniciativa: objetivo estratégico, hipótese de impacto, risco associado, métricas de acompanhamento e critério de continuidade ou interrupção.

Nessa lógica, os OKRs deixam de ser apenas metas declarativas e passam a orientar decisões de alocação de recursos. Já os OrKRs adicionam visibilidade sobre riscos críticos (tecnológicos, regulatórios, operacionais e de mercado), permitindo ação antecipada. O resultado esperado é maior disciplina de priorização e menor dispersão de esforço.

Como fechamento desta seção, o framework integrado posiciona o Grupo 3 para decidir com rapidez, aprender com evidências e preservar coerência entre estratégia, operação e desempenho econômico.

\section{Diferencial competitivo do Grupo 3}

O diferencial do Grupo 3 está na combinação de visão sistêmica e execução integrada. A proposta conecta decisões estratégicas, desenho de ecossistema, capacidade de engenharia e sustentabilidade financeira em uma arquitetura única de gestão.

\subsection{Diferencial pedagógico do modelo completo}

O modelo proposto fortalece a qualidade acadêmica do trabalho ao articular diferentes perspectivas de empreendedorismo em uma trilha única de decisão. Em vez de tratar os conteúdos de forma fragmentada, a abordagem integra formulação estratégica, validação prática e reflexão crítica sobre resultados.

Do ponto de vista formativo, a estrutura favorece o desenvolvimento simultâneo de competências analíticas, gerenciais e colaborativas. A autonomia discente é preservada na construção das hipóteses e na condução dos experimentos, enquanto a simulação de condições reais de mercado eleva a consistência das decisões apresentadas.

Outro ganho relevante é a combinação entre inovação incremental (aperfeiçoamento contínuo) e inovação de maior ruptura (novas propostas de valor), o que aumenta a robustez do raciocínio empreendedor e a aderência do projeto às demandas contemporâneas de negócio.

\subsection{Missão, entregáveis e foco do Grupo 3}

No arranjo geral do trabalho, o Grupo 3 tem como \textbf{missão} estruturar um negócio viável, convertendo ideias em um modelo com lógica econômica, coerência operacional e potencial de escala.

Para isso, os entregáveis são orientados por evidências e clareza de decisão: (i) \textit{Business Model Canvas} completo e internamente consistente; e (ii) modelagem objetiva de receitas, custos e parcerias-chave, com justificativas de priorização.

Esse foco amplia a qualidade técnica do capítulo ao explicitar a cadeia de valor do empreendimento, reforçando pensamento estratégico, estruturação organizacional e integração entre recursos, custos e monetização.

\subsection{Duplas funcionais e missões estratégicas}

Para viabilizar a execução do framework com alto grau de especialização e coordenação, o Grupo 3 organiza sua operação em três duplas complementares:

\begin{itemize}
    \item \textbf{Dupla A --- Arquitetos de Confiança (Hustler 1 + Hipster 1)}: tem como missão traduzir infraestrutura técnica em segurança como valor, criando um ecossistema no qual parceiros possam construir aplicações seguras sobre a base tecnológica proposta.
    \item \textbf{Dupla B --- Engenheiros de Defesa (Hacker 1 + Hipster 2)}: é responsável por produtizar scripts de defesa em módulos escaláveis de infraestrutura, assegurando também que a experiência de uso do \textit{dashboard} de segurança minimize fadiga de alertas e aumente a efetividade operacional.
    \item \textbf{Dupla C --- Guardiões do Caixa e Compliance (Hustler 2 + Hacker 2)}: atua na rastreabilidade das horas de PD\&I aplicadas a novos protocolos de segurança para viabilizar abatimento tributário de até 34\% via Lei do Bem, além de automatizar a geração de evidências técnicas exigidas pelo MCTI.
\end{itemize}

Como síntese gerencial, a Tabela~\ref{tab:trl-decisao-portfolio} consolida critérios de decisão por faixa de maturidade tecnológica, enquanto o Quadro~\ref{qua:governanca-integrada} articula a governança integrada entre tecnologia, finanças e conformidade.

Em termos de diferenciação competitiva, a proposta combina clareza estratégica, capacidade de coordenação e disciplina de governança, elevando a qualidade das decisões e a credibilidade do modelo perante avaliadores, parceiros e investidores.

\begin{table}[htb]
\centering
\caption{Matriz de maturidade e decisão para gestão do portfólio tecnológico.}
\label{tab:trl-decisao-portfolio}
\begin{tabular}{|p{0.15\textwidth}|p{0.35\textwidth}|p{0.38\textwidth}|}
\hline
\textbf{Faixa TRL} & \textbf{Foco principal} & \textbf{Decisão de gestão recomendada} \\
\hline
TRL 1--3 & Descoberta técnica e prova de conceito & Priorizar validação de hipótese crítica e reduzir escopo de incerteza. \\
\hline
TRL 4--6 & Prototipação e validação em ambiente relevante & Expandir testes, estruturar captação e preparar documentação de conformidade. \\
\hline
TRL 7--9 & Operação em ambiente real e escala & Acelerar comercialização, eficiência operacional e defesa competitiva. \\
\hline
\end{tabular}
\end{table}

\begin{table}[htb]
\centering
\caption{Quadro -- Governança para integração entre tecnologia, finanças e incentivos fiscais.}
\label{qua:governanca-integrada}
\begin{tabular}{|p{0.27\textwidth}|p{0.30\textwidth}|p{0.30\textwidth}|}
\hline
\textbf{Eixo de governança} & \textbf{Prática operacional} & \textbf{Valor agregado esperado} \\
\hline
Maturidade tecnológica & Revisão periódica de TRL/MRL com critério de avanço por evidência. & Menor risco de investimento prematuro e maior previsibilidade de entrega. \\
\hline
Engenharia financeira & Monitoramento de ROI/VPL por cenário e etapa de maturidade. & Melhor priorização de recursos e decisões de continuidade. \\
\hline
Incentivos fiscais & Rastreabilidade de horas, artefatos e resultados de PD\&I para Lei do Bem. & Economia tributária, segurança jurídica e reforço de caixa. \\
\hline
Conformidade regulatória & Registro de requisitos normativos e critérios de aceite por fase. & Redução de retrabalho e aceleração da entrada em mercado. \\
\hline
\end{tabular}
\end{table}

\section{Roteiro de implementação}

Nesta versão, o roteiro é tratado como \textbf{implementação executável} para uma startup de \textit{Hardware as a Service} (HaaS) em segurança da informação, com atuação inicial no Sul de Minas (Alfenas, Santa Rita do Sapucaí e Itajubá), priorizando agro e PMEs.

\subsection{Escopo implementado e tese de negócio}

O produto implementado é o \textbf{Kit Agro-Secure V1}: servidor de borda (\textit{home lab}), gateway de segurança, totem NFC e proteção elétrica, com monitoramento remoto de KPIs de continuidade operacional e segurança.

A tese de valor é objetiva: \textit{mesmo com internet instável e risco cibernético crescente, o cliente mantém operação crítica ativa, com rastreabilidade e governança de segurança em padrão corporativo}.

No posicionamento regional, a proposta aproveita a complementaridade do Sul de Minas: densidade do agronegócio (demanda operacional crítica), polos de engenharia eletrônica e telecomunicações (oferta técnica) e proximidade logística para atendimento presencial com SLA agressivo.

\subsection{Distribuição de responsabilidades entre as três duplas}

A incorporação integral do modelo Agro-Secure exige governança explícita de execução. Para reduzir sobreposição e aumentar accountability, as três duplas operam em frentes complementares, conforme a Tabela~\ref{tab:responsabilidades-duplas}.

\begin{table}[htb]
\centering
\caption{Distribuição de responsabilidades executivas entre as três duplas do Grupo 3.}
\label{tab:responsabilidades-duplas}
\begin{tabular}{|p{0.22\textwidth}|p{0.23\textwidth}|p{0.23\textwidth}|p{0.23\textwidth}|}
\hline
\textbf{Frente de execução} & \textbf{Dupla A (Arquitetos de Confiança)} & \textbf{Dupla B (Engenheiros de Defesa)} & \textbf{Dupla C (Caixa e Compliance)} \\
\hline
Arquitetura da oferta HaaS & Define padrão do kit, proposta técnica por segmento e desenho de valor percebido. & Especifica requisitos de firmware e operação segura do totem/edge. & Valida viabilidade econômica do padrão e impacto em CAPEX por cliente. \\
\hline
Instalação e operação em campo & Conduz diagnóstico inicial e desenho da topologia local. & Executa provisionamento, hardening e testes de failover no cliente. & Controla custo por instalação, deslocamento e produtividade por equipe. \\
\hline
SOC e monitoramento de KPIs & Define narrativa executiva dos relatórios para gestão do cliente. & Implementa coleta técnica (ameaças, uptime, temperatura, energia, link). & Consolida KPIs financeiros (MRR, payback, churn, margem de contribuição). \\
\hline
Validação comercial e escala & Estrutura pitch consultivo e parceiros regionais (cooperativas/associações). & Garante replicabilidade técnica dos kits e padrões de onboarding. & Modela cenários de expansão e necessidades de capital de giro. \\
\hline
Lei do Bem e conformidade & Enquadra hipóteses tecnológicas com foco em risco técnico real. & Produz evidências de experimentação, falhas e iterações técnicas. & Mantém timesheet, dossiê PD\&I e trilha para migração futura ao Lucro Real. \\
\hline
\end{tabular}
\end{table}

\subsection{Business Model Canvas completo}

\begin{table}[htb]
\centering
\caption{Business Model Canvas validável para startup InfoSec HaaS no Sul de Minas.}
\label{tab:bmc-haas-infosec}
\begin{tabular}{|p{0.30\textwidth}|p{0.62\textwidth}|}
\hline
\textbf{Bloco do BMC} & \textbf{Definição implementada} \\
\hline
Segmentos de clientes & Produtores rurais de café especial e grãos, cooperativas agrícolas e MMEs regionais do agro (laticínios, silos, comércio e logística). \\
\hline
Proposta de valor & Segurança \textit{enterprise} em formato \textit{plug and play}: servidores, roteadores seguros, totens NFC e monitoramento 24/7 sem necessidade de compra do hardware pelo cliente. \\
\hline
Canais & Venda consultiva in-loco, parcerias com cooperativas (ex.: Minasul e Cooxupé), feiras agropecuárias regionais e rede de indicação local. \\
\hline
Relacionamento com clientes & SLA agressivo, suporte presencial, onboarding operacional em campo e relatórios mensais executivos de ameaças bloqueadas, uptime preservado e estabilidade do ambiente. \\
\hline
Fontes de receita & Setup/instalação, mensalidade HaaS (hardware + firmware embarcado) e serviços de R\&D sob demanda para integrações e customizações críticas. \\
\hline
Recursos-chave & Estoque de hardware (servidor edge, gateway, NFC e energia), engenheiros de hardware/redes, SOC remoto em Alfenas e governança contábil para Lei do Bem. \\
\hline
Atividades-chave & Montagem e provisionamento de kits, instalação em campo (cabeamento/fixação), monitoramento de rede, manutenção e PD\&I em firmware e sincronização assíncrona segura. \\
\hline
Parcerias-chave & Fornecedores OEM (Brasil/Ásia), ISPs regionais, cooperativas agro, parceiros institucionais locais e escritório contábil com domínio em Lucro Real/Lei do Bem. \\
\hline
Estrutura de custos & CAPEX e depreciação de hardware, deslocamento/logística, folha técnica especializada, operação SOC 24/7, reposição de ativos e dispêndios de PD\&I. \\
\hline
\end{tabular}
\end{table}

\subsection{Funil de vendas e conversão regional}

Com base no anexo Agro-Secure, o funil comercial combina confiança local, prova técnica em campo e fechamento por evidência de valor econômico.

\begin{figure}[htb]
\centering
\fbox{\begin{minipage}{0.93\textwidth}
\textbf{Funil de vendas Agro-Secure (eixo Alfenas--Machado--Varginha)}\\
\textbf{Topo do funil (Dupla A):} prospecção por CNAE + introdução via cooperativas/sindicatos/associações.\\
\textbf{Meio do funil (Dupla B):} diagnóstico técnico e piloto de 30 dias com baseline de risco.\\
\textbf{Fundo do funil (Dupla A + Dupla C):} proposta HaaS com KPI executivo (uptime, ameaças bloqueadas, estabilidade) e contrato anual.\\
\textbf{Pós-venda e expansão (Dupla C + Dupla B):} renovação, upsell de integrações e escala por clusters regionais.
\end{minipage}}
\caption{Figura sugestiva do funil de vendas consultivo para conversão de pilotos em contratos HaaS.}
\label{fig:funil-vendas-haas}
\end{figure}

\begin{figure}[htb]
\centering
\fbox{\begin{minipage}{0.93\textwidth}
\textbf{Funil de metas de conversão (referência anual)}\\
\textbf{Etapa 1:} 15--20 prospects qualificados por território e CNAE.\\
\textbf{Etapa 2:} 3 pilotos ativos (30 dias) com coleta de KPIs operacionais.\\
\textbf{Etapa 3:} conversão comercial dos pilotos para contratos recorrentes.\\
\textbf{Etapa 4:} expansão via cooperativas (realista: 20 clientes no ano 1; otimista: 50).\\
\textbf{Meta financeira de referência:} MRR mês 12 entre R\$ 16.000,00 e R\$ 45.000,00.
\end{minipage}}
\caption{Figura sugestiva de funil de metas comerciais e financeiras para validação regional.}
\label{fig:funil-metas-conversao}
\end{figure}

\subsection{Arquitetura da solução e diagramas de validação}

As Figuras~\ref{fig:mapa-mental-haas} e~\ref{fig:usecase-campo} sintetizam, respectivamente, a arquitetura de valor e o fluxo operacional mínimo validado em campo.

\begin{figure}[htb]
\centering
\fbox{\begin{minipage}{0.92\textwidth}
\textbf{Mapa mental --- Oferta Agro-Secure V1}\\
\textbf{Núcleo:} Infraestrutura Segura HaaS\\
\textbf{Perímetro:} Gateway com firewall embarcado + controle de acesso NFC.\\
\textbf{Processamento local:} Servidor de borda para operação offline e sincronização assíncrona.\\
\textbf{Interface física:} Totem/Kiosk robusto para uso operacional em ambiente agro.\\
\textbf{Operação contínua:} SOC remoto com monitoramento de ameaças, energia, temperatura e disponibilidade.
\end{minipage}}
\caption{Figura sugestiva de mapa mental para arquitetura de oferta de valor.}
\label{fig:mapa-mental-haas}
\end{figure}

\begin{figure}[htb]
\centering
\fbox{\begin{minipage}{0.92\textwidth}
\textbf{Diagrama de caso de uso --- validação em campo}\\
Gestor da Fazenda $\rightarrow$ autentica via NFC no Totem Seguro.\\
Totem $\rightarrow$ valida acesso no Servidor Local HaaS.\\
Servidor Local $\rightarrow$ criptografa e sincroniza com nuvem/cooperativa.\\
SOC da startup $\dashrightarrow$ monitora anomalias, aplica atualização de firmware e registra KPIs.
\end{minipage}}
\caption{Figura sugestiva de caso de uso para operação local e monitoramento remoto.}
\label{fig:usecase-campo}
\end{figure}

\subsection{Backlog de validação e histórias de usuário}

O backlog prioriza \textbf{validação de disposição a pagar}, \textbf{adoção operacional} e \textbf{resiliência do kit} antes de qualquer escala de engenharia avançada.

\begin{table}[htb]
\centering
\caption{Backlog inicial orientado à validação de negócio (Epics e histórias).}
\label{tab:backlog-validacao}
\begin{tabular}{|p{0.20\textwidth}|p{0.20\textwidth}|p{0.52\textwidth}|}
\hline
\textbf{Épico} & \textbf{História} & \textbf{Critério de validação} \\
\hline
Disposição a pagar & HU 1.1 & Instalar 3 kits piloto em Alfenas por 30 dias e demonstrar ataques bloqueados, indisponibilidades evitadas e percepção de valor para conversão comercial. \\
\hline
Disposição a pagar & HU 1.2 & Entregar painel executivo mobile com métricas simples (ameaças bloqueadas, horas de operação preservadas e estabilidade da rede). \\
\hline
Adoção operacional & HU 2.1 & Prototipar totem com NFC e validar uso real por operadores de campo (luvas, poeira, rotina de turno), com taxa mínima de sucesso de autenticação. \\
\hline
Resiliência local & HU 3.1 & Garantir persistência local de eventos críticos (notas, pesagens, acessos) durante queda de link, com sincronização segura após retorno. \\
\hline
Conformidade PD\&I & HU 4.1 & Marcar tarefas inovadoras com etiqueta [PD\&I], registrar horas e evidências técnicas para formação contínua do dossiê Lei do Bem. \\
\hline
\end{tabular}
\end{table}

\subsection{Modelo técnico-operacional e unit economics do kit}

\begin{table}[htb]
\centering
\caption{Composição técnica e CAPEX médio do Kit Agro-Secure V1 por cliente.}
\label{tab:kit-capex-v1}
\begin{tabular}{|p{0.30\textwidth}|p{0.38\textwidth}|p{0.24\textwidth}|}
\hline
\textbf{Componente} & \textbf{Especificação de referência} & \textbf{Custo médio (R\$)} \\
\hline
Servidor Edge & Mini PC industrial fanless (16 GB RAM, SSD 500 GB) & 2.100,00 \\
\hline
Gateway de borda & Roteador corporativo (MikroTik/Ubiquiti) com VPN & 650,00 \\
\hline
Totem NFC & SBC + tela touch + leitor NFC + case reforçada & 1.500,00 \\
\hline
Energia & Nobreak senoidal puro 1200 VA & 1.200,00 \\
\hline
Miscelânea & Cabeamento, conectores, caixa de proteção e instalação & 450,00 \\
\hline
\textbf{CAPEX total por cliente} & \textbf{Kit físico alocado em comodato operacional} & \textbf{5.900,00} \\
\hline
\end{tabular}
\end{table}

\begin{table}[htb]
\centering
\caption{Precificação, payback e contribuição marginal do modelo HaaS.}
\label{tab:precificacao-payback}
\begin{tabular}{|p{0.46\textwidth}|p{0.34\textwidth}|}
\hline
\textbf{Indicador} & \textbf{Valor de referência} \\
\hline
Setup (instalação/auditoria) & R\$ 2.000,00 \\
\hline
Mensalidade HaaS & R\$ 850,00 \\
\hline
Receita no mês 1 por novo cliente & R\$ 2.850,00 \\
\hline
Exposição inicial de caixa por cliente & Aproximadamente R\$ 3.050,00 \\
\hline
Payback do hardware & Entre meses 5 e 6 \\
\hline
\end{tabular}
\end{table}

\subsection{Cenários de viabilidade e metas de validação}

\begin{table}[htb]
\centering
\caption{Cenários de viabilidade no primeiro ciclo anual de operação.}
\label{tab:cenarios-viabilidade}
\begin{tabular}{|p{0.25\textwidth}|p{0.20\textwidth}|p{0.20\textwidth}|p{0.20\textwidth}|}
\hline
\textbf{Métrica} & \textbf{Pessimista} & \textbf{Realista} & \textbf{Otimista} \\
\hline
Clientes ativos no ano 1 & 5 & 20 & 50 \\
\hline
CAPEX acumulado (R\$) & 20.000 & 80.000 & 180.000 \\
\hline
MRR no mês 12 (R\$) & 4.000 & 16.000 & 45.000 \\
\hline
Leitura executiva & Pivotar escopo para preservar caixa & Crescimento orgânico com estabilização & Escalar com capital externo \\
\hline
\end{tabular}
\end{table}

\subsection{Tratamento estatístico bayesiano da viabilidade}

Para integrar incerteza de mercado, execução e risco de capital intensivo, adota-se um modelo bayesiano de atualização de probabilidade de sucesso do negócio. Considera-se sucesso como atingir tração comercial com manutenção de continuidade operacional e evolução para equilíbrio econômico no horizonte de 12 a 14 meses.

\begin{table}[htb]
\centering
\caption{Evidências usadas na atualização bayesiana da probabilidade de sucesso.}
\label{tab:bayes-evidencias}
\begin{tabular}{|p{0.29\textwidth}|p{0.20\textwidth}|p{0.20\textwidth}|p{0.21\textwidth}|}
\hline
\textbf{Evidência observada} & \textbf{P(E$\mid$Sucesso)} & \textbf{P(E$\mid$Fracasso)} & \textbf{Razão de verossimilhança} \\
\hline
Demanda agro regional aquecida & 0,80 & 0,40 & 2,00 \\
\hline
Canal cooperativista estruturado & 0,75 & 0,45 & 1,67 \\
\hline
Dor técnica elevada (conectividade/ataques) & 0,85 & 0,50 & 1,70 \\
\hline
Restrição de crédito rural (evidência adversa) & 0,40 & 0,70 & 0,57 \\
\hline
Capacidade interna de execução (3 duplas) & 0,70 & 0,50 & 1,40 \\
\hline
\end{tabular}
\end{table}

Com probabilidade \textit{a priori} de 50\%, a atualização com as evidências da Tabela~\ref{tab:bayes-evidencias} resulta em probabilidade \textit{a posteriori} aproximada de \textbf{81,9\%} para sucesso do modelo, mantendo disciplina de risco operacional e financeiro.

\begin{table}[htb]
\centering
\caption{Atualização bayesiana dos cenários e MRR esperado no mês 12.}
\label{tab:bayes-cenarios-mrr}
\begin{tabular}{|p{0.20\textwidth}|p{0.16\textwidth}|p{0.18\textwidth}|p{0.18\textwidth}|p{0.16\textwidth}|}
\hline
\textbf{Cenário} & \textbf{Priors} & \textbf{P(E$\mid$Cenário)} & \textbf{Posterior} & \textbf{MRR mês 12 (R\$)} \\
\hline
Pessimista & 0,30 & 0,40 & 0,203 & 4.000 \\
\hline
Realista & 0,50 & 0,70 & 0,593 & 16.000 \\
\hline
Otimista & 0,20 & 0,60 & 0,203 & 45.000 \\
\hline
\multicolumn{4}{|r|}{\textbf{MRR esperado (valor esperado bayesiano)}} & \textbf{19.458} \\
\hline
\end{tabular}
\end{table}

Esse tratamento estatístico reforça a estratégia de execução por ondas: preserva prudência de caixa no curto prazo e, ao mesmo tempo, sustenta decisão de escala progressiva apoiada em evidências observáveis de conversão e retenção.

Como métrica de execução incremental, adota-se uma trilha em três ondas: (i) pilotos com 3 kits e conversão por valor demonstrado; (ii) expansão via parcerias com cooperativas; e (iii) consolidação regional com operação SOC 24/7 e backlog técnico escalável.

\begin{quadro}[htb]
\centering
\caption{Plano de execução por ondas com protagonismo das duplas.}
\label{qua:ondas-execucao-duplas}
\begin{tabular}{|p{0.20\textwidth}|p{0.20\textwidth}|p{0.16\textwidth}|p{0.14\textwidth}|p{0.20\textwidth}|}
\hline
\textbf{Onda} & \textbf{Meta de negócio} & \textbf{Dupla líder} & \textbf{Prazo de referência} & \textbf{Entregável crítico} \\
\hline
Onda 1 -- Pilotos & Validar conversão e disposição a pagar em campo & Dupla A & Meses 1--4 & 3 pilotos ativos, relatório executivo de valor e taxa de conversão monitorada \\
\hline
Onda 2 -- Escala regional & Acelerar aquisição com cooperativas e padronizar implantação & Dupla B & Meses 5--8 & Playbook técnico de instalação, onboarding e operação resiliente \\
\hline
Onda 3 -- Consolidação & Elevar MRR, governança e prontidão para expansão/captação & Dupla C & Meses 9--12 & Painel financeiro-operacional e dossiê PD\&I auditável \\
\hline
\end{tabular}
\end{quadro}

\begin{quadro}[htb]
\centering
\caption{Hipóteses críticas e indicadores de validação de negócio.}
\label{qua:hipoteses-validacao}
\begin{tabular}{|p{0.30\textwidth}|p{0.26\textwidth}|p{0.30\textwidth}|}
\hline
\textbf{Hipótese} & \textbf{Indicador-chave} & \textbf{Critério de sucesso} \\
\hline
Cliente paga pela continuidade operacional & Taxa de conversão piloto$\rightarrow$contrato & Conversão mínima de 40\% em até 45 dias \\
\hline
Kit reduz risco operacional real & Horas de operação preservadas em quedas de link & Evidência mensal de indisponibilidade evitada \\
\hline
Valor de segurança é compreendido pela gestão & Relatório de ameaças bloqueadas aceito na renovação & Renovação no mês 2 superior a 80\% \\
\hline
Modelo sustenta expansão regional & Relação MRR/CAPEX por cluster geográfico & Ponto de equilíbrio operacional até mês 14 \\
\hline
\end{tabular}
\end{quadro}

\subsection{Lei do Bem como eixo de vantagem competitiva}

A Lei do Bem (Lei n\textordmasculine\ 11.196/2005) é tratada como mecanismo estratégico de médio prazo. O desenho operacional separa atividades rotineiras de instalação das atividades de \textbf{risco tecnológico}, que são as elegíveis para incentivo fiscal quando houver migração sustentável ao regime de Lucro Real.

\begin{table}[htb]
\centering
\caption{Estrutura mínima do dossiê técnico de PD\&I para Lei do Bem.}
\label{tab:dossie-lei-bem}
\begin{tabular}{|p{0.32\textwidth}|p{0.58\textwidth}|}
\hline
\textbf{Seção do dossiê} & \textbf{Implementação prática na startup} \\
\hline
Escopo do projeto & Definição de objetivo técnico, período, incerteza e critério de sucesso por projeto embarcado. \\
\hline
Risco tecnológico & Registro do estado da arte, barreira técnica e hipótese de superação em ambiente agro de alta latência. \\
\hline
Evidência experimental & Prototipagem, teste de campo, falhas registradas, iterações e decisão técnica por versão. \\
\hline
Timesheet técnico & Apontamento de horas por tarefa [PD\&I], separando engenharia inovadora de operação rotineira. \\
\hline
Dispêndios elegíveis & Consolidação de mão de obra técnica, materiais de protótipo e custos vinculados a experimentação. \\
\hline
\end{tabular}
\end{table}

Como fechamento executivo desta seção, a implementação proposta transforma o roteiro em \textbf{plano de execução validável}, com entrega concreta de valor ao cliente, disciplina financeira para um modelo intensivo em hardware e construção antecipada de conformidade para capturar benefícios da Lei do Bem no momento econômico adequado.

Em síntese estratégica, a integração entre as três duplas permite converter complexidade técnica em resultado de negócio: Dupla A acelera tração e posicionamento, Dupla B assegura excelência operacional em campo e Dupla C protege sustentabilidade financeira e vantagem fiscal de longo prazo.

\section{Avaliação}

Para tornar os critérios de desempenho transparentes e comparáveis entre equipes, adota-se a Tabela~\ref{tab:criterios-avaliacao}, com pontuação total de 10,0 pontos.

Para maximizar desempenho em todos os critérios, recomenda-se que cada seção entregue três evidências mínimas: problema de negócio claramente delimitado, decisão metodológica justificada e resultado verificável (métrica, validação ou aprendizado documentado). Essa lógica reduz subjetividade na correção e fortalece a consistência da apresentação final.

Como fechamento avaliativo, a estratégia recomendada é sustentar cada afirmação com evidência objetiva, explicitando nexo entre problema, método e resultado, de modo a maximizar clareza, rigor e impacto da apresentação.

\begin{table}[htb]
\centering
\caption{Matriz de atendimento dos critérios de avaliação com argumentação de contrapartida.}
\label{tab:criterios-avaliacao}
\begin{tabular}{|p{0.23\textwidth}|p{0.39\textwidth}|p{0.24\textwidth}|p{0.08\textwidth}|}
\hline
\textbf{Critério} & \textbf{Argumento de contrapartida (como o trabalho atende)} & \textbf{Evidência objetiva no documento} & \textbf{Pontos} \\
\hline
Clareza da proposta & O problema regional, o ICP por CNAE e a tese de valor HaaS foram definidos com recorte geográfico e econômico explícito. & BMC validável, matriz de prospecção, funil com etapas e metas de conversão. & 1,0 \\
\hline
Aplicação dos conceitos & Os referenciais da disciplina foram operacionalizados em artefatos de gestão, engenharia, finanças e governança de valor. & Framework de cinco dimensões, backlog, hipóteses, cenários e plano por ondas. & 1,5 \\
\hline
Inovação e criatividade & A proposta integra hardware, firmware embarcado, edge computing e Lei do Bem como diferencial técnico-fiscal. & Kit Agro-Secure V1, arquitetura de operação offline, trilha de PD\&I e dossiê técnico. & 2,0 \\
\hline
Qualidade dos entregáveis & O material apresenta coerência executiva entre estratégia, operação, modelagem financeira e conformidade regulatória. & Tabelas de CAPEX/payback, cenários de viabilidade, quadros de responsabilidades e governança. & 1,5 \\
\hline
Apresentação e participação & A narrativa está estruturada para defesa oral objetiva, com argumentos verificáveis e divisão clara de protagonismo entre duplas. & Matriz de responsabilidades, funil comercial revisado, indicadores de validação e sínteses executivas por seção. & 4,0 \\
\hline
\textbf{Total} &  &  & \textbf{10,0} \\
\hline
\end{tabular}
\end{table}